% Track p3 -- Which charge is the photon, and the residual choice.
% Verification: investigations/photon_selection/verify/p3_identification.py (claims p3-1 ... p3-22).
% New references: sections/bib_p3.tex.
\section{Which charge is the photon, and the residual choice}
\label{p3:sec:main}

The second-rung construction gauges the charge $Q=\tfrac12J_1$ and carries it, in the curved extension, by the normal $SO(2)$ connection
\cite[Thms.~26.7, 26.8]{ES17}. At the third rung that charge is $Q_A=\tfrac12q_A=\tfrac12E_2E_3$, and the matter block of \cite[Section~8]{ES17} carries a
second candidate, the electric charge $B_{\rm em}$. This section decides which of the two is the photon on the family $F_3$ of \cref{p1:sec:main}, and what is
left to choose once a gauged subalgebra has aligned the vacuum. The results are as follows. The charge operator of $Q_A$ is half the involution $U$ that cuts
out $P_*$, so on the block $Q_A$ equals $\tfrac12$ on every state: it is half the fermion number, and it is anomalous together with $\su2_L$ (fixed). It is not
electromagnetism on that matter, whatever the vacuum (obstructed). The electric charge leaves the four aligned vacua massless, has four charged complex
fluctuations with the electron's charge magnitude, and acts on the block with the Standard-Model charges and without anomaly. With $B_{\rm em}$ gauged, the
massless vector of the projector current is the electromagnetic photon of the block, and the normal $U(1)$ is a separate abelian field (fixed). Of the three
counts against the normal $U(1)$ in \cite[Section~31]{ES17} only the sign count changes status, and a fourth count, the charge on the block, decides.
After alignment by $\mathfrak h_{\rm SM}$ four vacua remain. A Klein group of symmetries of the projector sector permutes them simply transitively, so no
alignment energy, kinetic term or current separates them (fixed). The overlaps with $P_*$, with the lepton projector and with the chirality halves do
not separate them either (fixed). The overlap with the half $P_L$ of the block on which $\su2_L$ acts selects the pair $\{\Pi_A,\Pi_{-W}\}$, and the two
members of the pair coincide on the half of the module that contains the block (fixed). A single point needs a coupling that sees the other half, for
instance the overlap with $\Pfam$ (constructed). Identifiers p3-$n$ refer to the \texttt{PASS} lines of \texttt{verify/p3\_identification.py}
(\cref{p3:tab:claims}).

\paragraph{Conventions.} Definitions and notation are those of \cref{p1:sec:main}; in particular $F_3$, $\Pi_{\pm A}$, $\Pi_{\pm W}$, $Q_A$, $B_{\rm em}$, $t^{L,R}_i$,
$t_Y$, $H_C$, $\rho_+$ and $\mathrm{vol}_+$ are as there. For a compact generator $X$ commuting with $J$ the charge operator is $\hat X:=-JX$. It is symmetric, and
$X=J\hat X$ acts as $i\hat X$ in the complex structure $J$ \cite[Section~8]{ES17}. The inner product on $\mathfrak k$ is $\langle X,Y\rangle=-\Tr(XY)/512$. For it the unit
blades $e_ie_j$ have norm one, $|t^L_3|^2=\tfrac18$ and $|t_Y|^2=\tfrac5{24}$. Orthonormal bases are $2\sqrt2\,t^L_i$, $\sqrt{24/5}\,t_Y$ and an orthonormal basis of
$\su3_C$ for $\mathfrak h_{\rm SM}$, with couplings $g_2$, $g_Y$, $g_3$, and the unit blades of $\so4_W\oplus\so6_C$ for $\mathfrak h_{\rm PS}$, with one coupling
$g$. We write $U_A=\operatorname{span}(E_2,E_3)$, $A_0=\operatorname{span}(E_8,E_9)$, $U_{\rm em}=\operatorname{span}(E_7,E_{13})$, $W_0=\operatorname{span}(E_{14},E_{15})$
and $P_{L,R}=P_*\tfrac12(\id\mp\Gamma)$. For an oriented four-plane $P\subset V_+$ with oriented orthonormal frame $(f_1,\dots,f_4)$, written as a $4\times8$ matrix $F$ in
the generator basis of $V_+$, the Plücker coordinates are $p_I(P)=\det F_I$ for the $70$ four-subsets $I$ of that basis, and $e_I$ is the blade of $I$ in increasing
order. We write $p_A$ and $p_W$ for $I=A$ and $I=W$. The family projector is $\Pfam=\tfrac18\prod_i(\id+X_i)$ \cite[Section~7]{ES17}, and
$Q_v=P_L+\gamma^0P_L\gamma^0$ \cite[Prop.~21.6]{ES17}.

\begin{definition}[The lepton projector]
\label{p3:def:plep}
$P_{\rm lep}$ is the spectral projector of the colour charge $Q_c=-JH_C$ for the eigenvalue $-3$. With $s_k=J\,E_{9+k}E_{15+k}$, $k=1,2,3$, one has
$Q_c=s_1+s_2+s_3$ and $P_{\rm lep}=\tfrac18\prod_k(\id-s_k)$.
\end{definition}

On $P_*$ the eigenvalue $-3$ of $Q_c$ is the colour singlet \cite[Prop.~8.4]{ES17}. \Cref{p3:prop:restated} shows that $P_{\rm lep}$ reproduces the row of
\cite[Prop.~21.6]{ES17}, where the colour-singlet lepton projector is named but not written out.

\subsection{Two charges on the matter}
\label{p3:sec:charges}

\begin{proposition}[The construction's charge is half of $U$; fixed, and restated]
\label{p3:prop:charges}
\begin{enumerate}[label=(\roman*),leftmargin=2.2em]
\item $\hat Q_A=-JQ_A=\tfrac12U$ on $\R^{512}$, where $U=-Jq_A$ is the third involution of $P_*$. Since $UP_*=P_*$ and $q_A$ commutes with $X$, $Z$ and $U$,
\[
q_AP_*=JP_*,\qquad \hat Q_AP_*=\tfrac12P_* .
\]
The first identity is the statement $q_A|_{V_*}=J$ of \cite[Section~8]{ES17} (p3-1).
\item On $V_*=\operatorname{im}P_*$ the operators $\Gamma$, $\hat T^3_L=-Jt^L_3$, $\hat Y=-Jt_Y$, $\hat Q_{\rm em}=-JB_{\rm em}$, $\hat Q_A$ and $\BmL=-JH_C/3$ commute. Their
joint spectrum is that of \cref{p3:tab:charges} (p3-2).
\item On the Weyl block $\tfrac12(\id-\Gamma)$ of the spectator action, of rank $256$, $\hat Q_A$ takes the values $\pm\tfrac12$, each $128$ times. There
$\hat Q_{\rm em}$ takes the values $\pm1$ ($16$ times each), $\pm\tfrac23$ ($48$), $\pm\tfrac13$ ($48$) and $0$ ($32$), the charges of the two families and two mirror
families of \cite[Prop.~10.12]{ES17} (p3-3).
\end{enumerate}
\end{proposition}

\begin{table}[t]
\centering
\small
\begin{tabular}{@{}lcccccccc@{}}
\toprule
field & $\Gamma$ & $\hat T^3_L$ & $\hat Y$ & $\hat Q_{\rm em}$ & $\hat Q_A$ & $\BmL$ & real dim. & complex dim.\\
\midrule
$Q_L$ & $-1$ & $\pm\tfrac12$ & $\tfrac16$ & $\tfrac23,\ -\tfrac13$ & $\tfrac12$ & $\tfrac13$ & $12+12$ & $6+6$\\
$L_L$ & $-1$ & $\pm\tfrac12$ & $-\tfrac12$ & $0,\ -1$ & $\tfrac12$ & $-1$ & $4+4$ & $2+2$\\
$u_R$ & $+1$ & $0$ & $\tfrac23$ & $\tfrac23$ & $\tfrac12$ & $\tfrac13$ & $12$ & $6$\\
$d_R$ & $+1$ & $0$ & $-\tfrac13$ & $-\tfrac13$ & $\tfrac12$ & $\tfrac13$ & $12$ & $6$\\
$e_R$ & $+1$ & $0$ & $-1$ & $-1$ & $\tfrac12$ & $-1$ & $4$ & $2$\\
$\nu_R$ & $+1$ & $0$ & $0$ & $0$ & $\tfrac12$ & $-1$ & $4$ & $2$\\
\bottomrule
\end{tabular}
\caption{Joint eigenvalues on $V_*$ (p3-2). Complex dimensions refer to the complex structure $J$; one internal Weyl component has four real, or two complex,
dimensions. The columns $\hat Y$ and $\hat Q_{\rm em}$ restate \cite[Prop.~8.4]{ES17}. The column $\hat Q_A$ is constant.}
\label{p3:tab:charges}
\end{table}

\begin{proof}
(i) $\hat Q_A=-J\tfrac12q_A=\tfrac12U$. As index sets, $X\sim I\cup A\cup C$ with $I$ the interior, $Z\sim A\cup W$ and $U\sim\{8,9\}\cup W\cup C$. The bivector $q_A$ shares two
indices with $X$ and with $Z$ and none with $U$, so by equation~(3) of \cite{ES17} it commutes with all three. $X$, $Z$, $U$ commute, so $UP_*=P_*$, that is
$-Jq_AP_*=P_*$; multiplying by $J$ and using $J^2=-\id$ gives $q_AP_*=JP_*$. (ii) The operators are the charge operators of commuting compact elements, together with
$\Gamma$, which commutes with every internal bivector; p3-2 diagonalises them jointly on $V_*$. (iii) $U$ commutes with $\Gamma$ and is traceless on the Weyl block;
p3-3 computes both spectra there.
\end{proof}

\begin{corollary}[What the construction's photon couples to; fixed, restated and obstructed]
\label{p3:cor:notem}
\begin{enumerate}[label=(\roman*),leftmargin=2.2em]
\item A gauge field coupled through $Q_A$ couples to every Weyl fermion of the block with the same charge $\tfrac12$, on both chiralities. In the complex structure
$J$ each of the sixteen Weyl fermions has fermion number one \cite[Section~8]{ES17}. The coupling is therefore to $\tfrac12F$, $F=3B+L$, a universal vector charge
(fixed).
\item No real $\lambda$ gives $\hat Q_AP_*=\lambda\hat Q_{\rm em}P_*$, since $\hat Q_{\rm em}$ vanishes on the neutrinos and $\hat Q_A$ does not (p3-4). The statement
involves only $P_*$, so the construction's charge is not the electric charge on the block for any vacuum projector (obstructed).
\item $\hat Q_{\rm em}$ has the same spectrum on $P_L$ as on $P_R$, so its cubic, gravitational and mixed colour anomalies cancel between the chiralities, and
$\Tr(P_L\hat Q_{\rm em}(\hat T^3_L)^2)=0$ (fixed; the anomaly conditions are inherited \cite{BouchiatIliopoulosMeyer72,GeorgiGlashow72}). For $\hat Q_A$ the same
traces vanish except $\Tr(P_L\hat Q_A(\hat T^3_L)^2)=4$. This is the mixed $SU(2)_L^2\,U(1)_A$ anomaly of \cite[Prop.~8.5]{ES17} at $a=\tfrac12$, $b=0$ (restated;
p1-28, p3-5). With $P_*$ as the matter and $\su2_L$ gauged, $\uu1_A$ cannot be gauged (obstructed under these hypotheses).
\end{enumerate}
\end{corollary}

\begin{proof}
(i) and (ii) read off \cref{p3:tab:charges}. (iii) is p3-5. The anomaly of $\hat Q_A$ is that of fermion number in the Standard Model, which $SU(2)_L$ instantons
violate \cite{tHooft76}.
\end{proof}

The universality of $Q_A$ on matter is therefore not a property of the vacuum. It is the factor $\tfrac12(\id+U)$ of $P_*$: the block is cut out as an eigenspace of
the construction's charge (\cref{p3:prop:charges}(i)), and on the rest of the spectator Weyl block the same charge takes the value $-\tfrac12$
(\cref{p3:prop:charges}(iii)). At the second rung the coupled model and the vacuum share the charge $\tfrac12J_1$ \cite[Section~26]{ES17}. At the third rung that
shared charge is half the fermion number.

\subsection{The counts against the normal \texorpdfstring{$U(1)$}{U(1)}, and the dichotomy}
\label{p3:sec:counts}

\Cref{p3:tab:counts} takes the three counts of \cite[Section~31]{ES17} to the third rung and adds the count of \cref{p3:cor:notem}. The four aligned vacua
$\Pi_0\in\{\Pi_{\pm A},\Pi_{\pm W}\}$ commute with both $Q_A$ and $B_{\rm em}$ (\cref{p1:thm:vacuum}), so both charges can be compared at the same vacuum.

\begin{table}[t]
\centering
\footnotesize
\begin{tabular}{@{}>{\raggedright\arraybackslash}p{3.1cm}>{\raggedright\arraybackslash}p{4.2cm}>{\raggedright\arraybackslash}p{4.4cm}>{\raggedright\arraybackslash}p{2.6cm}@{}}
\toprule
count & normal $U(1)$, generator $Q_A$ & $\uu1_{\rm em}$, generator $B_{\rm em}$ & change\\
\midrule
(i) a spinless test particle cannot see it \cite[Prop.~23.4, Thm.~30.1]{ES17}
& holds: it concerns the metric charge of the normal plane; the field couples through the spinor sector, at $\Pi_0$ to four complex scalars $\phi_w$ and to the block
& not applicable, since $B_{\rm em}$ is not a metric charge; it couples through the spinor sector to four complex scalars $\phi_a$ and to the block, and a neutral spinless field does not see the ordinary photon either
& none in substance (restated; interpretive reading)\\
(ii) the adapted generator does not commute with the physical generators \cite[Section~25]{ES17}
& removed: $[Q_A,\Pi_0]=0$, $D\Pi_0=0$, and the adapted correction is compact and commutes with the interior gammas \cite[Thm.~26.8]{ES17}, \cref{p1:thm:curved}
& removed in the same way
& none at $n=3$; removed already at $n=2$ (fixed)\\
(iii) it gravitates with the wrong sign \cite[Cor.~23.6, Thm.~30.4]{ES17}
& holds for the Kaluza--Klein vector of the normal plane
& not covered: the Maxwell term is adopted, and a positive one exists since $\langle B_{\rm em},B_{\rm em}\rangle>0$; read as a Kaluza--Klein vector, the plane $U_{\rm em}$ has the sign of time, like the normal plane
& from forced to adopted (constructed); the Kaluza--Klein reading is interpretive\\
(iv) the charge on the block, \cref{p3:cor:notem}
& universal $\tfrac12$, anomalous with $\su2_L$
& the Standard-Model charges, anomaly free
& decides (obstructed for $Q_A$, fixed for $B_{\rm em}$)\\
\bottomrule
\end{tabular}
\caption{The counts against the normal $U(1)$ at the third rung (p3-4, p3-5). In row (iii), $E_2,E_3,E_7,E_{13}$ square to $+\id$, as the time generator $\gamma^0=E_1$
does, and the colour generators to $-\id$ (p3-4).}
\label{p3:tab:counts}
\end{table}

Row (iii) rests on \cite[Cor.~23.6, Thm.~30.4]{ES17}: there a vector obtained by reduction on a plane of the sign of time has a kinetic term of the wrong sign
\cite{Hull98}. That theorem is about a metric origin of the field. An internal gauge field with an adopted Maxwell term is not covered by it. For a compact generator
the invariant form is positive (\cite[Thm.~2.3]{ES17}, \cref{p3:tab:counts}), so a Maxwell term of the healthy sign exists, as in ordinary Yang--Mills theory. If $\uu1_{\rm em}$
were itself obtained by reduction, its weak part would rotate the plane $U_{\rm em}$, of the sign of time, and the same mechanism would apply to that part and not to the
colour part $\tfrac16H_C$ (interpretive). The same holds for $Q_A$, whose plane $U_A$ also has the sign of time. Count (iii) therefore separates the two fields only by
their origin, and not by their generator.

The dichotomy at the end of \cite[Section~25]{ES17} reads: either the normal $U(1)$ is not electromagnetic, or the paired-mode projector is not the observation projector.
\cite[Section~26]{ES17} resolves it in one direction: if the normal $U(1)$ is to be electromagnetism, the paired-mode projector must go. At the third rung count (iv)
decides the first alternative for every projector, because it is a property of $P_*$ through its factor $U$: on the block matter the normal $U(1)$ is not electromagnetic
(obstructed). The replacement of the paired-mode projector is still available, as $\Pi_A$ or any point of the aligned set, and the charge that uses it is $B_{\rm em}$.
At $n=3$ both alternatives of the dichotomy hold, and neither is a choice.

\subsection{Identification at the aligned vacua}
\label{p3:sec:ident}

\begin{proposition}[Joint fixed points of the two charges; fixed]
\label{p3:prop:jointfixed}
The points of $F_3$ commuting with both $Q_A$ and $B_{\rm em}$ form six components: the isolated points $\Pi_{\pm(U_A\oplus U_{\rm em})}$, the isolated points $\Pi_{\pm R}$ with
$R=\operatorname{span}(E_8,E_9,E_{14},E_{15})$, and two copies of $\widetilde{\mathrm{Gr}}_2(\R^4)$, of dimension $4$,
\[
\{\Pi_{U_A\oplus P'}:P'\in\widetilde{\mathrm{Gr}}_2(R)\},\qquad\{\Pi_{U_{\rm em}\oplus P'}:P'\in\widetilde{\mathrm{Gr}}_2(R)\}.
\]
$\Pi_{\pm A}$ lie on the first of these, with tangent $A_0\otimes W_0$, and $\Pi_{\pm W}$ on the second (p3-6).
\end{proposition}

\begin{proof}
By \cref{p1:thm:fixed}, $\Pi_P$ commutes with $Q_A$ iff $P\supset U_A$ or $P\perp U_A$, and with $B_{\rm em}$ iff $P\supset U_{\rm em}$ or $P\perp U_{\rm em}$. The four
combinations give $P=U_A\oplus U_{\rm em}$; $P=U_A\oplus P'$ with $P'\subset(U_A\oplus U_{\rm em})^\perp=R$; $P=U_{\rm em}\oplus P'$; and $P=R$. An orientation of a direct sum
is a pair of orientations of the summands up to a common reversal. This gives $\widetilde{\mathrm{Gr}}_2(R)$ in the middle cases and two points in the outer ones. Finally
$A=U_A\oplus A_0$ and $W=U_{\rm em}\oplus W_0$ with $A_0,W_0\subset R$. p3-6 computes the joint tangent kernels $0$, $0$, $4$, $4$ at representatives.
\end{proof}

\begin{proposition}[Charges of the fluctuations; fixed]
\label{p3:prop:torus}
At each of $\Pi_{\pm A}$, $\Pi_{\pm W}$ the sixteen tangent directions decompose under the torus $\uu1_A\times\uu1_{\rm em}$ into complex lines of weights
\[
U_A\otimes U_{\rm em}:\ (1,1)\oplus(1,-1),\qquad U_A\otimes W_0:\ 2\times(1,0),\qquad A_0\otimes U_{\rm em}:\ 2\times(0,1),
\]
and four real neutral directions $A_0\otimes W_0$ (p3-7). Eight real directions are charged under each $\uu1$, four under both and four under neither, in agreement
with \cref{p1:thm:tangent}(iii),(iv).
\end{proposition}

\begin{proof}
$\rho_+(Q_A)$ is the rotation of $U_A$ at unit rate and $\rho_+(B_{\rm em})$ the rotation of $U_{\rm em}$; on $A\otimes W$ the first acts on the first factor and the
second on the second factor.
\end{proof}

\begin{theorem}[The photon of the aligned vacua; fixed, with the gauged generator a choice]
\label{p3:thm:photon}
Let $\Pi_0\in\{\Pi_{\pm A},\Pi_{\pm W}\}$. Let the internal connection carry a field $A_\mu$ along a compact generator $Q\in\mathfrak k$, acting on spinors by $Q$ and on the
projector by $\ad Q$, and adopt a Maxwell term.
\begin{enumerate}[label=(\alph*),leftmargin=2.2em]
\item For $Q=B_{\rm em}$ the four requirements on an electromagnetic photon of the block hold. (1) The generator is the electric charge: $\hat Q_{\rm em}=\hat T^3_L+\hat Y$ on
$V_*$, with the values of \cref{p3:tab:charges}. (2) $[B_{\rm em},\Pi_0]=0$, so the kinetic term gives no $A_\mu A^\mu$ term. (3) Charged fluctuations exist: the four complex
scalars $\phi_a$ of \cref{p1:thm:tangent}(iii), with current $j^\mu_{\Pi,\rm em}=128f^2\sum_a\Im(\bar\phi_aD^\mu\phi_a)$ (\cref{p1:thm:kinetic}). (4) The matter carries the
Standard-Model charges without anomaly (\cref{p3:cor:notem}(iii)). Both currents have the same unit of charge: $B_{\rm em}^2=-\hat Q_{\rm em}^2$ on $V_*$, with value $-1$ on
the electron, and $(\ad B_{\rm em})^2=-1$ on the charged tangent directions (p3-8). The scalars of the projector current carry the electron's charge magnitude, and no
normalisation is free.
\item For $Q=Q_A$ requirements (2) and (3) hold and (1) and (4) fail (\cref{p3:cor:notem}). The massless vector couples to $\tfrac12F$ on the block and to four complex
scalars $\phi_w$ of charge $1$, twice the fermion charge (p3-8). It is not the electromagnetic photon (obstructed).
\item $Q_A$ and $B_{\rm em}$ commute and are orthogonal. $Q_A$ commutes with $\mathfrak h_{\rm SM}$ and $\mathfrak h_{\rm PS}$ and lies in the centraliser $\so4_A$ of
$\mathfrak h_{\rm PS}$ in $\mathfrak k$ (p3-4). The normal $SO(2)$ connection of \cite[Thm.~26.8]{ES17}, whose plane is the surplus pair $E_2,E_3$ of the first
constituent, is the component of the connection along $Q_A$. It is a second abelian field, not the photon (fixed).
\item All twelve generators of $\mathfrak h_{\rm SM}$ commute with $\Pi_0$, so with $\mathfrak h_{\rm SM}$ gauged no gauge boson is massive at $\Pi_0$, and $\uu1_{\rm em}$
is singled out only by an electroweak stage. The projector's own neutral directions provide one. At $\Pi_\theta=\Ad(\exp\tfrac\theta2E_8E_{14})\Pi_A$, $\theta\notin\pi\Z$, the
mass matrix of $\mathfrak h_{\rm SM}$ has the eigenvalue $0$ nine times (gluons and photon), $m_W^2$ twice and $m_Z^2$ once, with
\[
\frac{m_W^2}{m_Z^2}=\frac1{1+\tfrac35\,g_Y^2/g_2^2},
\]
and its massless electroweak combination is $B_{\rm em}$ (p3-9). At $g_Y=g_2$ the ratio is $\tfrac58=1-\tfrac38$, the value $\cos^2\theta_W$ for the
$\sin^2\theta_W=\tfrac38$ of \cite[Thm.~21.1]{ES17}.
\end{enumerate}
\end{theorem}

\begin{proof}
(a) (1) is \cref{p3:prop:charges}(ii). (2) and (3) are \cref{p1:thm:vacuum,p1:thm:tangent,p1:thm:kinetic}. At $\Pi_{-A}=\id-\Pi_A$ the tangents change sign and nothing
else changes, and at $\Pi_{\pm W}$ \cref{p1:thm:tangent}(v) applies. For the unit: on $V_*$, $B_{\rm em}=J\hat Q_{\rm em}$ with $J$ commuting and $J^2=-\id$, so
$B_{\rm em}^2=-\hat Q_{\rm em}^2$; on the tangent, $[B_{\rm em},T_{a,7}]=T_{a,13}$ and $[B_{\rm em},T_{a,13}]=-T_{a,7}$ (\cref{p1:thm:tangent}). (b) p3-8 gives
$(\ad Q_A)^2=-1$ on $U_A\otimes W$ and $Q_A^2=-\tfrac14$ on $V_*$. (c) is p3-4 together with the form $\Omega=\Omega_{\rm Lorentz}+AQ$, $Q=\tfrac12ab$, of the curved
extension \cite[Thm.~26.8]{ES17}. (d) $\su3_C$ acts trivially on $V_+$. In the vector model of \cref{p1:lem:vector}, $p_\theta=p_{\operatorname{span}(E_2,E_3,E_9)}+uu^{\!\top}$ with
$u=\cos\theta\,E_8\pm\sin\theta\,E_{14}$. For $x\in\rho_+(\mathfrak h_{\rm SM})$, which kills $A$, $[x,uu^{\!\top}]=au^{\!\top}+ua^{\!\top}$ with $a=xu=\pm\sin\theta\,xE_{14}$, so
the mass matrix is $128\sin^2\theta$ times the Gram matrix of the vectors $g_a\rho_+(T_a)E_{14}$. The vectors $\rho_+(t^L_i)E_{14}$ are orthogonal of length $\tfrac12$,
because the $2\rho_+(t^L_i)$ are anticommuting complex structures, and $\rho_+(t_Y)E_{14}=-\rho_+(t^L_3)E_{14}$, because $B_{\rm em}$ kills $E_{14}$. In the basis of the
conventions the Gram matrix is $2g_2^2$ on the first two generators, and on $(2\sqrt2\,t^L_3,\sqrt{24/5}\,t_Y)$ it is
$\bigl(\begin{smallmatrix}2g_2^2&-\sqrt{12/5}\,g_2g_Y\\-\sqrt{12/5}\,g_2g_Y&\frac65g_Y^2\end{smallmatrix}\bigr)$, with eigenvalues $0$ and $2g_2^2+\tfrac65g_Y^2$. p3-9 computes
the full $12\times12$ matrix on $\R^{512}$.
\end{proof}

So at the third rung the photon of the aligned vacua is the $\uu1_{\rm em}$ gauge field of the internal connection, and the normal $U(1)$ is a separate abelian field. On
the block that field couples to fermion number, and with $\su2_L$ gauged it cannot be gauged on that matter (\cref{p3:cor:notem}). The construction's
identification of its photon with the normal connection is therefore not compatible with the matter at $n=3$ (obstructed). What the second rung established survives
with $B_{\rm em}$ in place of $Q_A$: the massless criterion of \cite[Cor.~26.2]{ES17} at a Lorentz-invariant vacuum, the charged-scalar current, and a curved extension in which
the vacuum is parallel (\cref{p1:thm:curved}). The theorem rests on the choice of the gauged generator, and $B_{\rm em}$ is fixed only once the Standard-Model embedding of
\cite[Section~8]{ES17} is adopted (constructed). The coupling of the second abelian field and the values of $f$ and of the gauge couplings are not decided (open).
\Cref{p3:thm:photon}(d) shows that the projector can itself carry an electroweak stage: its neutral doublet directions break $\mathfrak h_{\rm SM}$ to
$\su3_C\oplus\uu1_{\rm em}$ with the tree-level relation of one doublet \cite{Glashow61,Weinberg67}. Whether the vacuum moves there is decided by the potential, which
\cref{p3:sec:robust} examines.

\subsection{The residual choice}
\label{p3:sec:residual}

By \cref{p1:prop:K,p1:thm:fixed}, with $\mathfrak h_{\rm SM}$ or $\mathfrak h_{\rm PS}$ gauged with positive couplings, the alignment energy vanishes on $F_3$ exactly at
$\Pi_A$, $\Pi_{-A}$, $\Pi_W$ and $\Pi_{-W}$. This subsection determines what separates these four points.

\begin{lemma}[Overlaps are Plücker functions; fixed]
\label{p3:lem:plucker}
For every oriented four-plane $P\subset V_+$, $\mathrm{vol}_P=\sum_Ip_I(P)\,e_I$, and for every operator $M$ on $\R^{512}$
\[
\Tr(\Pi_PM)=\tfrac12\Tr M+\tfrac12\sum_Ip_I(P)\Tr(e_IM).
\]
The four-blade components $\Tr(e_IM)$ vanish for $M=P_*$, $P_{\rm lep}$ and $\tfrac12(\id\pm\Gamma)$. The nonzero ones are $\Tr(e_A\Pfam)=64$,
$\Tr(e_AP_L)=-\Tr(e_WP_L)=32$, their negatives for $P_R$, $\Tr(e_AQ_v)=-\Tr(e_WQ_v)=64$ and $\Tr(e_A\Pi_A)=256$ (p3-10).
\end{lemma}

\begin{proof}
The frame vectors are orthonormal, so their Clifford product equals their exterior product, $\sum_I\det F_I\,e_I$. The blades $e_I$ are symmetric involutions with
$\Tr(e_Ie_J)=512\,\delta_{IJ}$. For the components it suffices to list index sets. The eight products of $X$, $Z$, $U$ have the index sets $\emptyset$, $I\cup A\cup C$,
$A\cup W$, $\{8,9\}\cup W\cup C$ and their symmetric differences, none of which is a four-subset of $V_+$. With $\Gamma$, whose index set is all eighteen generators, the
products include $\Gamma X=\chi_W$ and $\Gamma XZ=-\chi_W\mathrm{vol}_+=-\mathrm{vol}_A$, so $P_L=\tfrac1{16}(\id+X)(\id+Z)(\id+U)(\id-\Gamma)$ contains
$\tfrac1{16}(\mathrm{vol}_A-\chi_W)$ and no other four-blade of $V_+$. The products of the $s_k$ have index sets of sizes $12$, $4$ inside $C$, and $8$ equal to $A\cup W$.
Among the products of the $X_i$ only $X_1=\mathrm{vol}_A$ is a four-blade of $V_+$ (\cref{p3:prop:restated}). $\gamma^0$ commutes with every $e_I$, which doubles the components
of $P_L$ in $Q_v$.
\end{proof}

\begin{proposition}[The overlaps; fixed]
\label{p3:prop:overlaps}
\begin{enumerate}[label=(\roman*),leftmargin=2.2em]
\item On all of $F_3$, $\Tr(\Pi P_*)=\Tr(\Pi P_{\rm lep})=32$ and $\Tr(\Pi\tfrac12(\id\pm\Gamma))=128$ are constant, while
\[
\Tr(\Pi_P\Pfam)=32(1+p_A)=\tfrac14\Tr(\Pi_P\Pi_A),\qquad \Tr(\Pi_PP_L)=16(1+p_A-p_W)=\tfrac12\Tr(\Pi_PQ_v),
\]
and $\Tr(\Pi_PP_R)=16(1-p_A+p_W)$ (p3-12).
\item $\Tr(\Pi P_*\Pi P_*)=\|P_*\Pi P_*\|^2\le\Tr(\Pi P_*)=32$, with equality iff $[\Pi,P_*]=0$. The points of $F_3$ commuting with $P_*$ are $\Pi_{\pm A}$ and $\Pi_{\pm W}$,
which are isolated, and the two components $\{\Pi_{U_A\oplus W'}\}$ and $\{\Pi_{A_0\oplus W'}\}$, $W'\in\widetilde{\mathrm{Gr}}_2(W)$, of dimension $4$ (p3-12, p3-13).
\item At the four points the linear and quadratic overlaps are those of \cref{p3:tab:overlaps} (p3-11).
\end{enumerate}
\end{proposition}

\begin{table}[t]
\centering
\small
\begin{tabular}{@{}lccccl@{}}
\toprule
$M$ & $\Pi_A$ & $\Pi_{-A}$ & $\Pi_W$ & $\Pi_{-W}$ & $\Tr(\Pi_PM)$ on $F_3$\\
\midrule
$P_*$ & $(32,32)$ & $(32,32)$ & $(32,32)$ & $(32,32)$ & $32$\\
$P_{\rm lep}$ & $(32,32)$ & $(32,32)$ & $(32,32)$ & $(32,32)$ & $32$\\
$\tfrac12(\id\pm\Gamma)$ & $(128,128)$ & $(128,128)$ & $(128,128)$ & $(128,128)$ & $128$\\
$P_L$ & $(32,32)$ & $(0,0)$ & $(0,0)$ & $(32,32)$ & $16(1+p_A-p_W)$\\
$P_R$ & $(0,0)$ & $(32,32)$ & $(32,32)$ & $(0,0)$ & $16(1-p_A+p_W)$\\
$Q_v$ & $(64,64)$ & $(0,0)$ & $(0,0)$ & $(64,64)$ & $32(1+p_A-p_W)$\\
$\Pfam$ & $(64,64)$ & $(0,0)$ & $(32,16)$ & $(32,16)$ & $32(1+p_A)$\\
\bottomrule
\end{tabular}
\caption{Overlaps $(\Tr\Pi M,\ \Tr\Pi M\Pi M)$ at the four aligned points, and the linear overlap on all of $F_3$ (p3-11, p3-12).}
\label{p3:tab:overlaps}
\end{table}

\begin{proof}
(i) is \cref{p3:lem:plucker} with $\Tr M$ and the components listed there. (ii) $P_*$ is a contraction, so $\|P_*\Pi P_*\|^2\le\|\Pi P_*\|^2=\Tr(P_*\Pi P_*)=32$, with equality
iff $(\id-P_*)\Pi P_*=0$, that is iff $\Pi P_*=P_*\Pi P_*$ is symmetric, that is iff $\Pi$ commutes with $P_*$. Every $\mathrm{vol}_P$ commutes with $J$, $\Gamma$ and
$\mathrm{vol}_+=-Z$, so $\Pi_P$ commutes with $P_*$ iff $\mathrm{vol}_P$ commutes with $\chi_W$ and with $q_A$. $\Ad(\chi_W)$ acts on $V_+$ as the reflection $-1$ on $W$, $+1$ on $A$.
It fixes $\mathrm{vol}_P$ iff $P=(P\cap A)\oplus(P\cap W)$ with $\dim(P\cap W)$ even. With \cref{p1:thm:fixed} for $q_A$ this leaves $P=\pm A$, $P=\pm W$, and $P=A'\oplus W'$ with
$A'\in\{U_A,A_0\}$ and $W'$ a two-plane in $W$. p3-13 checks that the four points are isolated and that the components have dimension $4$. (iii) is the computation
p3-11.
\end{proof}

\begin{lemma}[A comass inequality; fixed]
\label{p3:lem:comass}
For every oriented four-plane $P\subset V_+=A\oplus W$, $|p_A(P)|+|p_W(P)|\le1$, with equality iff $P\in\{\pm A,\pm W\}$.
\end{lemma}

\begin{proof}
Let $M_A=(\langle a_i,f_j\rangle)$ and $M_W=(\langle w_i,f_j\rangle)$ for an oriented orthonormal frame $(f_j)$ of $P$, so that $M_A^{\!\top}M_A+M_W^{\!\top}M_W=\id$. The two
symmetric matrices commute, so in a common eigenbasis they are $\operatorname{diag}(\cos^2\theta_i)$ and $\operatorname{diag}(\sin^2\theta_i)$, with
$\theta_i\in[0,\tfrac\pi2]$ the principal angles between $P$ and $A$. Hence $|p_A|=\prod_i\cos\theta_i$ and $|p_W|=\prod_i\sin\theta_i$. For any two indices $i\neq j$,
\[
|p_A|+|p_W|\le\cos\theta_i\cos\theta_j+\sin\theta_i\sin\theta_j=\cos(\theta_i-\theta_j)\le1 .
\]
Equality for every pair forces all $\theta_i$ to equal one $\theta$, and then $\cos^4\theta+\sin^4\theta=1-2\cos^2\theta\sin^2\theta=1$ forces $\theta\in\{0,\tfrac\pi2\}$,
that is $P=\pm A$ or $P=\pm W$. p3-14 checks the identities at $300$ random planes and finds the maximum $1$ of $p_A-p_W$ only at $A$ and $-W$.
\end{proof}

\begin{corollary}[The extremes of the separating overlaps; fixed]
\label{p3:cor:extremes}
On $F_3$, $0\le\Tr(\Pi P_L)\le32$, with the maximum attained exactly at $\Pi_A$ and $\Pi_{-W}$ and the minimum exactly at $\Pi_{-A}$ and $\Pi_W$. Likewise
$0\le\Tr(\Pi\Pfam)\le64$, with the maximum only at $\Pi_A$ and the minimum only at $\Pi_{-A}$.
\end{corollary}

\begin{proof}
By \cref{p3:prop:overlaps}(i) and \cref{p3:lem:comass}, $|p_A-p_W|\le1$ with $p_A-p_W=1$ iff $P\in\{A,-W\}$ and $=-1$ iff $P\in\{-A,W\}$. Also $|p_A|\le1$ with equality
iff $P=\pm A$, since $|p_A|=\prod\cos\theta_i$.
\end{proof}

\begin{proposition}[A Klein group of the projector sector; fixed]
\label{p3:prop:klein}
Let $\mathsf C(\Pi)=\id-\Pi$, let $\mathsf D(\Pi)=\tfrac12(\id+\mathrm{vol}_+(2\Pi-\id))$ be the duality of \cref{p1:prop:duality}, and let $\mathsf E=\mathsf C\mathsf D$, so
that $\mathsf E(\Pi)=\tfrac12(\id+\mathrm{vol}_+)-\mathrm{vol}_+\Pi$.
\begin{enumerate}[label=(\roman*),leftmargin=2.2em]
\item $\{1,\mathsf C,\mathsf D,\mathsf E\}$ is a Klein four-group of $K$-equivariant involutions of $F_3$, with $\mathsf E(\Pi_P)=\Pi_{-P^\perp}$. It acts simply transitively on
the four points: $\mathsf D$ exchanges $\Pi_A$ and $\Pi_W$, $\mathsf C$ exchanges $\Pi_A$ and $\Pi_{-A}$, and $\mathsf E$ exchanges $\Pi_A$ and $\Pi_{-W}$ (p3-16).
\item Each element preserves every mass matrix $\Tr([T,\Pi][T',\Pi])$, $T,T'\in\mathfrak k$, and, for connections valued in $\mathfrak k$ and the interior Lorentz algebra, the
kinetic term and the current $j_\Pi$ of every $Q\in\mathfrak k$ (p3-16).
\item $\mathsf E$ is the identity on the half $Z=+1$, the range of $\tfrac12(\id+Z)$, which contains $P_*$: $\mathsf E(\Pi)\tfrac12(\id+Z)=\Pi\tfrac12(\id+Z)$, and
$P_*\mathsf E(\Pi)P_*=P_*\Pi P_*$. On that half $\Pi_A=\Pi_{-W}=\tfrac12(\id-\chi_W)$ and $\Pi_W=\Pi_{-A}=\tfrac12(\id+\chi_W)$, and there $t^L_i=\Pi_A\tau_i$ and
$t^R_i=\Pi_W\tau_i$. On the half $Z=-1$, $\Pi_A=\Pi_W$ and $\Pi_{-A}=\Pi_{-W}$ (p3-17).
\end{enumerate}
Hence no alignment energy, kinetic term or current separates any two of the four points. No coupling that depends on $\Pi$ only through $\Pi\tfrac12(\id+Z)$ separates
$\Pi_A$ from $\Pi_{-W}$, or $\Pi_{-A}$ from $\Pi_W$.
\end{proposition}

\begin{proof}
$\mathrm{vol}_+$ is $K$-invariant, commutes with the interior gammas and with every $\Pi_P$, and squares to $\id$ (\cref{p1:lem:census,p1:prop:duality}). $\mathsf C$ and
$\mathsf D$ commute, since both $\mathsf C\mathsf D(\Pi)$ and $\mathsf D\mathsf C(\Pi)$ equal $\tfrac12(\id-\mathrm{vol}_+(2\Pi-\id))$. (ii) For $T$ in $\mathfrak k$ or in the
interior Lorentz algebra, $[T,\mathsf C\Pi]=-[T,\Pi]$, $[T,\mathsf D\Pi]=\mathrm{vol}_+[T,\Pi]$ and $[T,\mathsf E\Pi]=-\mathrm{vol}_+[T,\Pi]$. Since $\mathrm{vol}_+$ is orthogonal and
commutes with $[T,\Pi]$, the traces are preserved. $\mathrm{vol}_+$ is parallel for such connections, so $D_\mu$ transforms in the same way. (iii) On $Z=+1$,
$\mathrm{vol}_+=-\id$, so $\mathsf E(\Pi)=\Pi$ there. $\mathrm{vol}_A=\mathrm{vol}_+\chi_W$ equals $-\chi_W$ there and $+\chi_W$ on $Z=-1$; with
$t^L_i=\tfrac12(\id-\chi_W)\tau_i$ \cite[Thm.~8.2]{ES17} the remaining identities follow.
\end{proof}

On the half of the module that contains the matter, the vacuum $\Pi_A$ is the support projector of $\su2_L$, and $\Pi_W$ that of $\su2_R$. The residual choice is
therefore, on the matter, the choice of which weak $\su2$ acts inside the image of the vacuum. It is not a choice of spacetime chirality: the chirality halves have
constant overlap, and the vector-like $Q_v$ separates exactly as the chiral $P_L$ does (\cref{p3:tab:overlaps}).

\begin{proposition}[Normalisers; fixed]
\label{p3:prop:normaliser}
\begin{enumerate}[label=(\roman*),leftmargin=2.2em]
\item The fixed vectors of $\rho_+(\mathfrak h_{\rm SM})$ and of $\rho_+(\mathfrak h_{\rm PS})$ in $V_+$ are exactly $A$. A reflection of $W$ conjugates $\su2_L$ onto $\su2_R$.
Hence every element of $K$ normalising $\mathfrak h_{\rm SM}$ fixes each of $\Pi_A$, $\Pi_{-A}$, $\Pi_W$, $\Pi_{-W}$ (p3-18).
\item $g=E_9E_{15}=\exp(\tfrac\pi2E_9E_{15})$ normalises $\mathfrak h_{\rm PS}$, exchanges $\su2_L$ and $\su2_R$, and sends $\mathrm{vol}_A\mapsto-\mathrm{vol}_A$,
$\mathrm{vol}_W\mapsto-\mathrm{vol}_W$ and $X\mapsto-X$. So the normaliser of $\mathfrak h_{\rm PS}$ identifies $\Pi_A$ with $\Pi_{-A}$ and $\Pi_W$ with $\Pi_{-W}$, by
elements that do not preserve $P_*$. No element of either normaliser maps $\Pi_{\pm A}$ to $\Pi_{\pm W}$ (p3-18).
\item $g_{AW}=\prod_i\tfrac1{\sqrt2}(\id+a_iw_i)$, with $(a_i)=(E_2,E_3,E_8,E_9)$ and $(w_i)=(E_7,E_{13},E_{14},E_{15})$, is an element of $\mathrm{Spin}(8)$ with
$\Ad(g_{AW})\mathrm{vol}_A=\mathrm{vol}_W$ and $\Ad(g_{AW})\mathrm{vol}_W=\mathrm{vol}_A$, which normalises neither $\mathfrak h_{\rm SM}$ nor $\mathfrak h_{\rm PS}$. Moreover
$\mathrm{vol}_A\mathrm{vol}_W=\mathrm{vol}_+$ (p3-18).
\end{enumerate}
\end{proposition}

\begin{proof}
(i) $\Ad(k)$ acts on $F_3$ through the rotation $R=\rho(k)$ of $V_+$ (\cref{p1:thm:classify}). If $\Ad(k)$ normalises $\mathfrak h$, $R$ preserves the fixed space $A$ of
$\rho_+(\mathfrak h)$ and its complement $W$. For $\mathfrak h_{\rm SM}$, $R|_W$ normalises the derived algebra $\su2_L$ of $\rho_+(\mathfrak h_{\rm SM})$. A rotation of $W$ of
determinant $-1$ is a rotation times a reflection, and it maps $\su2_L$ onto $\su2_R$; so $\det R|_W=1$, hence $\det R|_A=1$, and $R$ fixes the oriented planes $A$ and $W$. (ii)
$\Ad(g)$ is the rotation by $\pi$ in the plane $(E_9,E_{15})$, of determinant $-1$ on $A$ and on $W$; it preserves $W$ and $V_-$, hence $\mathfrak h_{\rm PS}$. $\Gamma$ commutes
with $g$ and $\Ad(g)\chi_W=-\chi_W$. The last statement follows from (i), which applies to $\mathfrak h_{\rm PS}$ as well. (iii) $\Ad(\exp\tfrac\pi4a_iw_i)$ rotates the plane
$(a_i,w_i)$ by $\tfrac\pi2$; the signs and the failure to normalise are computed in p3-18.
\end{proof}

With $\mathfrak h_{\rm SM}$ gauged the four points are therefore distinct points of the gauged theory, although the projector sector cannot tell them apart. With
$\mathfrak h_{\rm PS}$ gauged, $\Pi_A$ and $\Pi_{-A}$ are related by a left--right exchange of the Pati--Salam algebra, which the matter block breaks.

\subsection{Robustness}
\label{p3:sec:robust}

\begin{lemma}[Invariant couplings keep fixed critical points; inherited, with proof]
\label{p3:lem:robust}
Let $H\subset K$ be connected with Lie algebra $\mathfrak h$, let $f$ be a smooth $H$-invariant function on $F_3$, and let $\Pi_0\in F_3$ be fixed by $H$. Then $df_{\Pi_0}$ is
an $H$-invariant covector. If $T_{\Pi_0}F_3$ contains no nonzero $\mathfrak h$-invariant vector, $\Pi_0$ is a critical point of $f$. In particular $\Pi_0$ stays critical
for $V_0+tf$ for every $t$ whenever $V_0$ is $H$-invariant too, and it stays a nondegenerate local minimum as long as the Hessian of $V_0+tf$, an $H$-invariant
symmetric form, stays positive definite.
\end{lemma}

\begin{proof}
For $h\in H$, $f\circ h=f$ and $h\Pi_0=\Pi_0$ give $df_{\Pi_0}\circ dh_{\Pi_0}=df_{\Pi_0}$. The metric $\Tr(d\Pi\,d\Pi)$ is $K$-invariant (\cref{p1:lem:vector}), so the gradient
is an $H$-invariant tangent vector, and for a connected group the invariant vectors are those killed by $\ad\mathfrak h$ \cite{Michel71,Palais79}. The Hessian at a critical
point is invariant for the same reason.
\end{proof}

\begin{proposition}[Invariant tangent vectors; fixed]
\label{p3:prop:invariants}
The dimensions of the $\mathfrak h$-invariant subspaces of the tangent spaces at $\Pi_A$ and $\Pi_W$, and so at $\Pi_{-A}$ and $\Pi_{-W}$, whose tangent spaces are the same
up to sign, are: $\mathfrak h_{\rm SM}$, $\mathfrak h_{\rm PS}$, $\su2_L$ and $\uu1_Y\oplus\su3_C$: $0$; $\uu1_{\rm em}\oplus\su3_C$: $8$, the neutral components $A\otimes W_0$;
$\uu1_A$: $8$; $\uu1_A\oplus\uu1_{\rm em}$: $4$ (p3-19).
\end{proposition}

\begin{corollary}[Matter couplings that respect $\mathfrak h$; fixed, with an inherited principle]
\label{p3:cor:matter}
Let the vacuum energy on $F_3$ be the sum of an alignment energy and of couplings to matter that respect $\mathfrak h\in\{\mathfrak h_{\rm SM},\mathfrak h_{\rm PS}\}$: functions of $\Pi$
and of operators fixed by $\mathfrak h$, such as the overlaps of \cref{p3:prop:overlaps}, Yukawa-type kernels commuting with $\mathfrak h$, or the fermion loop of a Dirac operator
with $\mathfrak h$-invariant couplings \cite[Section~21]{ES17}. Such an energy is $H$-invariant, since $f(\Ad(h)\Pi;M)=f(\Pi;\Ad(h^{-1})M)$, and $\Pi_{\pm A}$, $\Pi_{\pm W}$ remain
critical points of it. All linear overlaps of \cref{p3:tab:overlaps}, including the one with $\Pfam$, which does not commute with $\su3_C$, are $\mathfrak h_{\rm SM}$-invariant
functions on $F_3$, and their first variations vanish at the four points along all of $\mathfrak k$ (p3-20). With only $\uu1_{\rm em}\oplus\su3_C$ gauged the eight neutral
directions are invariant, and an invariant coupling can move the vacuum along the $\uu1_{\rm em}$-fixed component $\widetilde{\mathrm{Gr}}_4(A\oplus W_0)$ of
\cref{p1:thm:fixed}.
\end{corollary}

\begin{proof}
\Cref{p3:lem:robust} with \cref{p3:prop:invariants}. The invariance of the overlaps is p3-20; for $\Pfam$ it also follows from $\Tr(\Pi_P\Pfam)=\tfrac14\Tr(\Pi_P\Pi_A)$ and
$[\mathfrak h_{\rm SM},\Pi_A]=0$.
\end{proof}

What invariant couplings can do is change the Hessian. Write $\Pi(x)=\Ad(e^X)\Pi_0$ with $X=\sum x_{aw}\tfrac12aw$, as in \cref{p1:thm:tangent}, and
\[
c_{\rm SM}=6g_2^2+\tfrac65g_Y^2,\qquad c_{\rm PS}=12g^2 .
\]

\begin{proposition}[Hessians at the four points; fixed]
\label{p3:prop:hessian}
At each of the four points $V_{\mathfrak h_{\rm SM}}(\Pi(x))=128\,c_{\rm SM}|x|^2+O(x^3)$ and $V_{\mathfrak h_{\rm PS}}(\Pi(x))=128\,c_{\rm PS}|x|^2+O(x^3)$. The overlaps expand as
\[
\Tr(\Pi(x)P_L)=\begin{cases}32-8|x|^2+O(x^3)&\text{at }\Pi_A,\ \Pi_{-W},\\ 8|x|^2+O(x^3)&\text{at }\Pi_{-A},\ \Pi_W,\end{cases}\qquad
\Tr(\Pi(x)\Pfam)=\begin{cases}64-16|x|^2+\cdots&\text{at }\Pi_A,\\ 16|x|^2+\cdots&\text{at }\Pi_{-A},\\ 32+O(x^4)&\text{at }\Pi_{\pm W}\end{cases}
\]
(p3-21). For $V=\kappa V_{\mathfrak h_{\rm SM}}-\lambda\Tr(\Pi P_L)-\mu\Tr(\Pi\Pfam)$ every Hessian at the four points is a multiple of the identity. $\Pi_A$ and $\Pi_{-W}$ are
nondegenerate local minima for all $\kappa\ge0$, $\lambda,\mu\ge0$ not all zero. $\Pi_W$ is a local minimum iff $\lambda<16\kappa c_{\rm SM}$, and $\Pi_{-A}$ iff
$\lambda+2\mu<16\kappa c_{\rm SM}$. Beyond these values all sixteen directions become unstable at once.
\end{proposition}

\begin{proof}
Since $[T_a,\Pi_0]=0$, $\|[T_a,\Pi(x)]\|^2=\|[[T_a,X],\Pi_0]\|^2+O(x^3)$. $[T_a,X]$ corresponds to $x\,\rho_+(T_a)^{\!\top}$ on $A\otimes W$, so the sum is
$128\,\tr(x\,C\,x^{\!\top})$ with $C=\sum_ag_a^2\rho_+(T_a)\rho_+(T_a)^{\!\top}$ on $W$. For $\mathfrak h_{\rm SM}$ the Casimir of $\su2_L$ and $\rho_+(t_Y)^2=-\tfrac14$ give
$C=(6g_2^2+\tfrac65g_Y^2)\id_W$, and $\su3_C$ gives nothing; for $\mathfrak h_{\rm PS}$ each of the four directions of $W$ lies in three unit blades, so $C=12g^2\id_W$. The
overlaps follow from \cref{p3:prop:overlaps}(i) with $p_A(\Pi(x))=1-\tfrac12|x|^2+O(x^4)$ near $A$, and $p_W=\det(O(x))=O(x^4)$ there; near $\pm W$ the roles are exchanged.
p3-21 checks the second derivatives $\Tr([X,[X,\Pi_0]]P)$ exactly and by finite differences. The stability conditions compare $128\kappa c_{\rm SM}$ with $8\lambda$, and with
$8\lambda+16\mu$ at $\Pi_{-A}$.
\end{proof}

The sign of $\lambda$ is a choice. For $\lambda<0$ the roles of $\{\Pi_A,\Pi_{-W}\}$ and $\{\Pi_{-A},\Pi_W\}$ are exchanged. The Hessians found here are the multiples of
the identity in the sixteen-dimensional space of $\mathfrak h_{\rm SM}$-invariant forms of \cref{p1:cor:critical}; a coupling with a component of the form
$N\otimes\omega_Y$ would split the doublet directions.

\subsection{The complete selection}
\label{p3:sec:selection}

The alignment result used is \cref{p1:prop:K,p1:thm:fixed}: for positive couplings $V_{\mathfrak h_{\rm SM}}\ge0$ on $F_3$, with zero set exactly $\{\Pi_{\pm A},\Pi_{\pm W}\}$.

\begin{theorem}[Selection on $F_3$]
\label{p3:thm:selection}
Assume:
\begin{enumerate}[label=(H\arabic*),leftmargin=2.6em]
\item the vacuum projector is a point of $F_3$ (constructed: \cref{p1:sec:answer});
\item $\mathfrak h_{\rm SM}$ is gauged, embedded as in \cite[Section~8]{ES17}, with couplings $g_3,g_2,g_Y>0$ (constructed);
\item the sign hypothesis: the gauge fluctuations contribute $\kappa V_{\mathfrak h_{\rm SM}}$ to the vacuum energy with $\kappa\ge0$ (inherited mechanism of vacuum alignment
\cite{Peskin80,Preskill81}; its sign is proved for vector-like theories \cite{Witten83} and assumed here);
\item the vacuum couples to the block by $-\lambda\Tr(\Pi P_L)$ with $\lambda>0$ (constructed, sign chosen);
\item optionally, it couples to the family block by $-\mu\Tr(\Pi\Pfam)$ with $\mu>0$ (constructed).
\end{enumerate}
Let $V=\kappa V_{\mathfrak h_{\rm SM}}-\lambda\Tr(\Pi P_L)-\mu\Tr(\Pi\Pfam)$, with $\mu=0$ when (H5) is not assumed.
\begin{enumerate}[label=(\alph*),leftmargin=2.2em]
\item Under (H1)--(H3) the minimum set is $\{\Pi_{\pm A},\Pi_{\pm W}\}$. Overlaps with $P_*$, $P_{\rm lep}$ and $\tfrac12(\id\pm\Gamma)$, linear or quadratic, take equal values at
these points and do not separate them.
\item Under (H1)--(H4) the global minimum value is $-32\lambda$, attained exactly at $\Pi_A$ and $\Pi_{-W}$, for every $\kappa\ge0$. These two points coincide on the half
$Z=+1$ that contains $P_*$, and they are exchanged by $\mathsf E$, which preserves every term of $V$ except the one of (H5). $\Pi_W$ and $\Pi_{-A}$ remain critical, and are
local minima iff $\lambda<16\kappa c_{\rm SM}$.
\item Under (H1)--(H5) the global minimum is $\Pi_A$ alone, with value $-32\lambda-64\mu$.
\item At the selected vacuum all of $\mathfrak h_{\rm SM}$ is unbroken, the massless vector of \cref{p3:thm:photon}(a) is the photon of $\uu1_{\rm em}$, and the field along $Q_A$ is
a second abelian field.
\end{enumerate}
\end{theorem}

\begin{proof}
(a) is \cref{p1:prop:K,p1:thm:fixed} and \cref{p3:tab:overlaps}. (b) $V\ge\kappa\cdot0-32\lambda$ by \cref{p3:cor:extremes}, with equality iff $V_{\mathfrak h_{\rm SM}}=0$ and
$\Tr(\Pi P_L)=32$, that is iff $\Pi\in\{\Pi_A,\Pi_{-W}\}$, where both conditions hold. The coincidence and the role of $\mathsf E$ are \cref{p3:prop:klein}; the overlap with $P_L$
depends on $\Pi$ only through $\Pi\tfrac12(\id+Z)$. Stability is \cref{p3:prop:hessian}. (c) With $\mu>0$ equality requires in addition $\Tr(\Pi\Pfam)=64$, which holds only at
$\Pi_A$ (\cref{p3:cor:extremes}). (d) is \cref{p1:thm:vacuum} and \cref{p3:thm:photon}. p3-22 minimises $V$ from random starts for several couplings, including
$\lambda>16\kappa c_{\rm SM}$ and $\kappa=0$, and finds exactly these minima.
\end{proof}

Two features of the theorem could have come out otherwise. The overlap with $P_*$ itself, which the matter block suggests first, is constant on $F_3$ and selects nothing,
and the selection is made by the weak-chirality half $P_L$. Once that overlap is present the alignment energy is not needed for the selection; it only stabilises the
other two critical points. The pair $\{\Pi_A,\Pi_{-W}\}$ is one projector on the half of the module that carries the matter, so for physics on the block the selected
vacuum is unique (interpretive). A unique point of $F_3$ needs a coupling that sees the other half. The overlap with $\Pfam$ is such a coupling, but $\Pfam$ carries no colour,
and adopting it amounts to choosing $\Pi_A$, since $\Tr(\Pi\Pfam)=\tfrac14\Tr(\Pi\Pi_A)$ on $F_3$ (constructed). The rank-$256$ projector $\Pi_A$ is the projector $P_1$ of the
projector family \cite[Thm.~7.5]{ES17}, on which colour and a chiral doublet first coexist \cite[Thms.~7.7, 7.9]{ES17}.

\begin{proposition}[Restated data of the projectors; restated]
\label{p3:prop:restated}
$X_1=J^{(1)}J^{(2)}=\mathrm{vol}_A$, so $P_1=\Pi_A$, with $\dim\mathfrak g_{\rm phys}(P_1)=51$, $\dim\mathfrak k(P_1)=27$ and a three-dimensional annihilator, as in
\cite[Thm.~7.5]{ES17}, and $\Pfam\le\Pi_A$. The projector $P_{\rm lep}$ of \cref{p3:def:plep} has rank $64$, commutes with the interior gammas and with $J$, and with $P_*$,
$\Pfam$ and $Q_v$ reproduces the alignment energies $3$, $5$, $9$, $5$ and stabilisers $37$, $23$, $9$, $25$ of \cite[Prop.~21.6]{ES17} (p3-15).
\end{proposition}

\begin{proof}
$X_1=e^{(1)}_2e^{(1)}_3e^{(2)}_2e^{(2)}_3=E_2E_3E_8E_9$. The other statements are the computation p3-15.
\end{proof}

\subsection{The answer}
\label{p3:sec:answer}

\begin{enumerate}[label=(\alph*),leftmargin=2.2em]
\item \emph{Which $\uu1$ is the photon.} At the third rung the photon of the aligned vacua is the $\uu1_{\rm em}$ gauge field of the internal connection, generated by
$B_{\rm em}=-\tfrac12E_7E_{13}+\tfrac16H_C$. It is massless at $\Pi_{\pm A}$ and $\Pi_{\pm W}$, its projector current is carried by four complex scalars of the electron's
charge magnitude, and on the block it couples with the Standard-Model charges without anomaly (fixed, given the adopted embedding, which is constructed).
\item \emph{The construction's charge.} $Q=\tfrac12J_1$ becomes $Q_A=\tfrac12E_2E_3$, whose charge operator is $\tfrac12U$. On the block it is half the fermion number, the
same on every state, and anomalous with $\su2_L$ (fixed; anomaly restated). It is not electromagnetism (obstructed). The normal $SO(2)$ connection that carries it is a
separate abelian field (fixed), and its coupling is not decided (open).
\item \emph{The counts and the dichotomy.} Of the three counts against the normal $U(1)$, the spinless count is unchanged, the adapted-generator count is removed for both
charges, and the sign count changes from forced to adopted; a fourth count, the charge on the block, decides. Both alternatives of the dichotomy of
\cite[Section~25]{ES17} hold at $n=3$ (fixed; the Kaluza--Klein reading of the sign count is interpretive).
\item \emph{Which vacuum.} Alignment by $\mathfrak h_{\rm SM}$ leaves $\Pi_{\pm A},\Pi_{\pm W}$ (fixed, \cref{p1:thm:fixed}). A Klein group of symmetries of the projector sector
permutes them simply transitively; no alignment energy, kinetic term or current separates them (fixed). The weak-chirality overlap $-\lambda\Tr(\Pi P_L)$, $\lambda>0$, selects
$\{\Pi_A,\Pi_{-W}\}$, which is one projector on the matter half (fixed under constructed hypotheses). A single point requires a coupling off that half, such as
$-\mu\Tr(\Pi\Pfam)$, which selects $\Pi_A$ (constructed).
\item \emph{Robustness.} $\mathfrak h_{\rm SM}$- or $\mathfrak h_{\rm PS}$-invariant couplings, including matter couplings, keep the four points critical, because their
tangent spaces carry no invariant vector (fixed; principle inherited). They can only change the Hessians, which for the couplings considered are multiples of the identity;
$\Pi_W$ and $\Pi_{-A}$ lose stability for $\lambda>16\kappa c_{\rm SM}$ and $\lambda+2\mu>16\kappa c_{\rm SM}$ (fixed).
\item \emph{Open.} The dynamical origin and sign of the overlap couplings, the scales $f$, $g$, $\kappa$, $\lambda$, $\mu$, the coupling of the second abelian field, whether the
half $Z=-1$ is physical, and the origin of the Maxwell term of $\uu1_{\rm em}$ are open.
\end{enumerate}

\begin{table}[t]
\centering
\footnotesize
\begin{tabular}{@{}>{\raggedright\arraybackslash}p{8.4cm}>{\raggedright\arraybackslash}p{4.3cm}>{\raggedright\arraybackslash}p{2.2cm}@{}}
\toprule
statement & class & checks\\
\midrule
\cref{p3:def:plep}, \cref{p3:prop:restated}: $P_{\rm lep}$, $P_1=\Pi_A$, the table of ES17 Prop.~21.6 & restated (ES17 Thm.~7.5, Prop.~21.6) & p3-15\\
\cref{p3:prop:charges}: $\hat Q_A=\tfrac12U$, $q_AP_*=JP_*$, charge table, Weyl block & fixed; $q_AP_*=JP_*$ and the $\hat Y$, $\hat Q_{\rm em}$ columns restated & p3-1 to p3-3\\
\cref{p3:cor:notem}: $Q_A$ is $\tfrac12F$ on the block, not the electric charge, anomalous; $\uu1_{\rm em}$ anomaly free & fixed; obstructed; anomaly restated; anomaly conditions inherited & p3-4, p3-5\\
\cref{p3:tab:counts}: the four counts and the dichotomy & restated (i); fixed (ii), (iv); constructed and interpretive (iii) & p3-4, p3-5\\
\cref{p3:prop:jointfixed}: joint fixed points of $Q_A$ and $B_{\rm em}$ & fixed & p3-6\\
\cref{p3:prop:torus}: torus weights of the fluctuations & fixed & p3-7\\
\cref{p3:thm:photon}: the photon is $\uu1_{\rm em}$; common charge unit; $Q_A$ field separate; $m_W^2/m_Z^2$ & fixed; gauged generator constructed; (b) obstructed & p3-4, p3-8, p3-9\\
\cref{p3:lem:plucker}, \cref{p3:prop:overlaps}: overlaps as Plücker functions, values, commutant of $P_*$ & fixed & p3-10 to p3-13\\
\cref{p3:lem:comass}, \cref{p3:cor:extremes}: comass inequality, extremes of the separating overlaps & fixed & p3-12, p3-14\\
\cref{p3:prop:klein}: Klein group, $\mathsf E$ trivial on the matter half & fixed & p3-16, p3-17\\
\cref{p3:prop:normaliser}: normalisers of $\mathfrak h_{\rm SM}$, $\mathfrak h_{\rm PS}$; $g_{AW}$ & fixed & p3-18\\
\cref{p3:lem:robust}: invariant couplings keep fixed critical points & inherited \cite{Michel71,Palais79}, proof given & ---\\
\cref{p3:prop:invariants}, \cref{p3:cor:matter}: invariant tangent vectors, matter couplings & fixed; invariance of induced energies inherited & p3-19, p3-20\\
\cref{p3:prop:hessian}: Hessians and stability conditions & fixed & p3-21\\
\cref{p3:thm:selection}: selection of $\{\Pi_A,\Pi_{-W}\}$, and of $\Pi_A$ with (H5) & fixed under (H1), (H2), (H4), (H5) constructed and (H3) inherited & p3-22\\
unique vacuum for physics on the block; Kaluza--Klein reading of the sign count & interpretive & ---\\
origin and sign of the overlap couplings; scales; the second abelian field & open & ---\\
\bottomrule
\end{tabular}
\caption{Claims of \cref{p3:sec:main} and their classes. The checks are the \texttt{PASS} lines printed by \texttt{verify/p3\_identification.py}.}
\label{p3:tab:claims}
\end{table}
